\paragraph{IV 21, \S1: invertible sheaves on projective bundles and their automorphisms (cote25, sites98--105).}
The section begins with the classification of invertible Modules on $P(\underline E)$ and develops it into a description of morphisms and automorphisms of projective bundles. Theorem1.1 and Corollaries1.3--1.5 describe the partition of the base by the integer twist and the remaining invertible Module on the base; Remark1.2 separates the rank-zero, rank-one and higher-rank parts. The proof first establishes uniqueness, then reduces existence to the field case and invokes III4.6.5 with the vanishing from III2; its preliminary constructions and the canonical projective-bundle isomorphisms refer to ChapterII. At site101 the author explicitly revises the exposition of the field case, asks to begin directly with the graded module of sections, and replaces the connected-base version by the formulation using locally constant integer-valued functions $\underline Z(S)$. The resulting map from the constant group scheme to the Picard scheme is assigned to the planned ChapterV, at sites101--102. Site102 makes a further revision: the author proposes stating1.1 without its initial rank restriction, separates the uniqueness assertions according to nonempty fibres and positive-dimensional fibres, and moves Remark1.2 into the proof. The same page introduces degree, and sites102--104 describe morphisms of a fixed degree through an invertible Module $\underline M$ and a homomorphism involving $\underline{\mathrm{Sym}}^n(\underline E)$. The cases of degree0 and degree1 are identified in the manuscript as the constant and the linear morphisms. Theorem1.6 and Corollaries1.7--1.8 treat isomorphisms and their local descriptions; Remark1.9 names the projective group and its absolute version $GP(r)$. Its handwritten instruction at site105 asks that the representability argument be carried out using a determinant open; the position of the possible dual accent in that marginal formula is above the first $E$, and its presence remains doubtful. Several revisions require the reader's attention without changing the text: site98 prints $\underline O_P(1)$ next to a final twist $(n)$ in the formula of1.1; site101 retains the printed $q$-operators in the proposed field-case exposition; and site102 says ``$\underline L$ inversible sur $S$'' in a reformulation whose preceding discussion concerns invertible Modules on $P$. The overprinted parenthesis/prime-like mark in the pullback formula at site103 remains uncertain. Large diagonal marks occur at sites98--103; their graphic form is recorded, but their function does not securely establish deletion of entire arguments. In particular, the unsupported paragraph-wide deletion wrappers at site100 have been removed, and the insertion connector at site101 is not treated as cancelling ``comme Module''. The page before this section, site97, contains a severely faint related text whose uncertain readings remain marked after inspection; it has supplied no reading for these clearer pages.
